\begin{afterword}
\summaryhead

The post follows ``Where Recursive Justification Hits Bottom,'' which concluded that it is fine to use induction and Occam's razor to assess induction and Occam's razor. A commenter had compared that view with W. W. Bartley's pancritical rationalism, and the author sets out what makes the view distinctive, without saying whether other philosophers share it.

The view comes from designing a self-modifying AI: if you would not want the AI to stop using induction, your philosophy should not call induction unjustifiable. The AI should treat its own parts like anything else in the world. Reflective consistency is a side effect of trying to win, not a goal. Jaynes's rule to use all available information licenses using induction on itself. One should inspect one's foundations in order to improve them, as simplicity was redefined so that Maxwell's equations count as simpler than Thor, while keeping ordinary results. ``Does induction work?'' may be answered by induction; ``Why does induction work?'' is still open. A loop of justification through the meta-level is not circular logic, though the author has not yet formalized the difference.

\responsehead

The post is candid about what it is. It says the author is ``far from the first person'' to consider reflection, reports the comparison with Bartley, and ends by saying the theory is not yet ``in hand.''

Much of it restates the previous post. Jaynes's rule, the claim that the point is to win rather than to be consistent, and the claim that reflective loops differ from circular logic all appear in ``Where Recursive Justification Hits Bottom.'' New here are the AI origin of the view, its naturalism, and the separation of ``does'' from ``why.''

The features have established names in philosophy, which the post does not give. Treating the study of one's own reasoning as one more natural inquiry is Quine's naturalized epistemology, which answered the charge of circularity in 1969. The difference between reflective loops and circular logic is close to the difference between rule-circular and premise-circular arguments, and Van Cleve and Papineau had argued that rule-circularity need not be vicious. The standard objection, that counterinduction can support itself in the same way, is the one the previous post answers with its anti-Laplacian minds. These precedents support the post, and the Stanford Encyclopedia entries cited in the notes give the arguments on both sides.

The comparison the post was prompted by turns on justification. Bartley gave up the demand that beliefs be justified; the previous post says that everything needs justification. The commenter pointed to this difference, and the post does not discuss it.

One step needs a premise the post does not state. The test for philosophies, that one which would not want an AI to stop using induction ``had better not label induction as unjustifiable,'' assumes that the label is advice to stop. Hume, the classic source of the label, held that inferences from experience come from custom, not reasoning, and that only ``a fool or madman'' would reject experience as a guide. The post does not say why a philosophy like his could not want the AI to keep induction.

The formal theory the post looks forward to was later attempted by the author and Marcello Herreshoff. They found that the most straightforward formalism of a self-modifying agent meets a ``L\"obian obstacle,'' which technical methods avoid while leaving ``the underlying puzzles of rational coherence'' only partly addressed. That bears on the post's suspicion that a correctly built AI would need to do nothing special.

In short: a candid list of features, several repeated from the previous post and others known in epistemology under other names, with one test for philosophies that assumes, without saying why, that calling induction unjustified is advice to stop using it.
\end{afterword}
